% Passes 5356--5361: canonical PSL2(q) pair space, Steinberg square,
% all-odd Hecke/Wedderburn formula, and q-mod-8 compatibility clock.
\subsection{Canonical two-line pair space and the all-odd Steinberg-square law}
\label{sec:pass5356-pair-steinberg}

The local fifteen-point fiber isolated at $q=5$ in Pass~5336 is not an
isolated representation-theoretic object.  For every odd prime power $q$, the
$q+1$ lines through a point of $W(3,q)$ carry the natural projective
$\PSL_2(q)$ action, and their unordered pairs form the canonical homogeneous
space
\[
 \Lambda_q=\binom{\mathbf P^1(q)}2
 \cong \PSL_2(q)/N(T_{\rm split}),
 \qquad
 |\Lambda_q|=\frac{q(q+1)}2,
 \qquad
 |N(T_{\rm split})|=q-1.
\]
At $q=5$, $\Lambda_5$ is exactly the $A_5/V_4$ fiber of Pass~5336.  For
general $q$ we use only the canonical incident-line pair space; we do
\emph{not} assume that an arbitrary minimum $K_0$ shell has this fiber.

Let $\mathrm{St}_q$ denote the $q$-dimensional complex Steinberg module of
$\PSL_2(q)$.  Since
\[
 \mathbb C[\mathbf P^1(q)]\cong \mathbf 1\oplus \mathrm{St}_q,
\]
and the symmetric square of the point-permutation module splits into its
diagonal monomials and its off-diagonal monomials, Pass~5356 gives the exact
characteristic-zero identity
\[
 \boxed{
 \mathbb C[\Lambda_q]\cong \operatorname{Sym}^2(\mathrm{St}_q).
 }
\]
The direct elementwise character check
\[
 \chi_{\Lambda_q}(g)
 =\frac{\chi_{\mathrm{St}}(g)^2+\chi_{\mathrm{St}}(g^2)}2
\]
was replayed for $q=3,5,7,11,13$.  In particular, the $q=5$ decomposition
\[
 15=1+4+2\cdot5
\]
is the first nontrivial member of an all-odd Steinberg-square family.

\paragraph{Orbital rank.}
Write $m=(q-1)/2$.  The pair stabilizer is a dihedral group of order $2m$.
A Burnside count of its orbits on $\Lambda_q$ gives
\[
 \boxed{
 r(q)=\dim\operatorname{End}_{\PSL_2(q)}\mathbb C[\Lambda_q]
 =\begin{cases}
 (3q+9)/4,&q\equiv1\pmod4,\\[2mm]
 (3q+7)/4,&q\equiv3\pmod4.
 \end{cases}}
\]
The identity contributes $|\Lambda_q|$ fixed pairs; each nonidentity split
rotation fixes the base pair, the unique rotational involution contributes an
additional $m$ pairs when $m$ is even, and each of the $m$ reflections fixes
$m+1$ pairs.  Thus the rank-six algebra at $q=5$ is not uniform in $q$.

\subsection{All-odd complex Wedderburn theorem for the pair Hecke algebra}
\label{sec:pass5360-pair-wedderburn}

The exact permutation-character decomposition of $\PSL_2(q)$ on two-subsets,
as computed by Darafsheh and Pournaki, shows that every irreducible constituent
occurs with multiplicity one or two.  More precisely, for $q\equiv1\pmod4$ the
multiplicity-two terms are the Steinberg character together with the relevant
principal-series characters indexed by $i\equiv0\pmod4$, while for
$q\equiv3\pmod4$ they are the discrete-series terms indexed by
$j\equiv2\pmod4$.  Counting them gives
\[
 d(q)=\left\lfloor\frac{q+3}{8}\right\rfloor.
\]
Consequently Pass~5360 upgrades the finite Pass~5359 intersection-tensor
census to the all-odd theorem
\[
 \boxed{
 \mathcal A_q:=\operatorname{End}_{\PSL_2(q)}\mathbb C[\Lambda_q]
 \cong
 \mathbb C^{\,s(q)}\oplus M_2(\mathbb C)^{\,d(q)},
 }
\]
where
\[
 s(q)=r(q)-4d(q).
\]
Hence
\[
 \boxed{
 \dim Z(\mathcal A_q)=r(q)-3d(q),
 \qquad
 \dim[\mathcal A_q,\mathcal A_q]=3d(q).
 }
\]
The direct intersection tensors at
$q=3,5,7,11,13,17,19,23$ agree exactly.  The first values are
\[
 \begin{array}{c|c|c}
 q&r(q)&\mathcal A_q\\\hline
 3&4&\mathbb C^4\\
 5&6&\mathbb C^2\oplus M_2(\mathbb C)\\
 7&7&\mathbb C^3\oplus M_2(\mathbb C)\\
 11&10&\mathbb C^6\oplus M_2(\mathbb C)\\
 13&12&\mathbb C^4\oplus M_2(\mathbb C)^2.
 \end{array}
\]
This is a theorem about the complex canonical pair space.  It does not identify
the characteristic-two footprint module.

\subsection{Characteristic-two firewall and the $q\bmod8$ Klein clock}
\label{sec:pass5361-qmod8-clock}

The characteristic-zero cancellation used above cannot simply be reduced
modulo two.  Put
\[
 V_2=\mathbb F_2[\mathbf P^1(q)],
 \qquad
 \mathbf j=\sum_x e_x,
 \qquad
 \varepsilon\!:\,V_2\to\mathbb F_2.
\]
Because $q+1$ is even,
\[
 \varepsilon(\mathbf j)=0.
\]
Transitivity gives
$\dim\operatorname{Hom}_{\PSL_2(q)}(V_2,\mathbf1)=1$, generated by
augmentation.  Therefore every equivariant map $V_2\to\mathbf1$ kills
$\mathbf j$, so no equivariant projection onto $\langle\mathbf j\rangle$
exists.  Thus the trivial submodule does not split off in characteristic two.
This is the explicit firewall preventing the Steinberg-square theorem from
being misused as a proof of the still-open all-odd binary footprint-rank law.

There is, however, an exact arithmetic compatibility.  The complex pair
character contains its exceptional conjugate pair precisely for
$q\equiv1,3\pmod8$, while the characteristic-two Weil pair from the rank-three
point module descends to $\mathbb F_2$ precisely for
$q\equiv1,7\pmod8$ (and otherwise requires $\mathbb F_4$).  These are the two
quadratic characters
\[
 \chi_{-2}(q),\qquad \chi_2(q),
\]
and
\[
 \chi_{-2}(q)\chi_2(q)=\chi_{-1}(q).
\]
Equivalently,
\[
 \boxed{
 (\mathbb Z/8\mathbb Z)^\times
 \longrightarrow \{\pm1\}^2,
 \qquad
 q\longmapsto(\chi_{-2}(q),\chi_2(q))
 }
\]
is the full Klein-four character clock.  The two bits come from independent
characteristic-zero and characteristic-two constructions; no isomorphism
between their exceptional constituents is asserted.

\paragraph{Evidence boundary.}
The all-odd statements proved in Passes~5356--5361 concern the canonical
$\PSL_2(q)$ pair space and its complex centralizer algebra, plus the elementary
characteristic-two nonsplitting statement and the arithmetic $q\bmod8$ clock.
The general equality
\[
 \operatorname{rank}_{2}F=\frac{q(q^2+1)}2
 \quad\Longleftrightarrow\quad
 \ker_2(F^T)=C_W
\]
remains open beyond the verified anchors.  Likewise the $q=11$ dual
weight-$20$ existence/distance problem and the Hoffman shortened-code minimum
distance remain unchanged.
